% =============================================================================
% Code Reviewer: an arXiv-style technical report of the whole codebase.
% Compile with:
%   tectonic main.tex
% =============================================================================
\documentclass[11pt,a4paper]{article}

% ---- packages -----------------------------------------------------------------
\usepackage[utf8]{inputenc}
\usepackage[T1]{fontenc}
\usepackage{lmodern}
\usepackage{microtype}
\usepackage[margin=1in]{geometry}
\usepackage{setspace}
\usepackage{titling}
\usepackage{titlesec}
\usepackage{authblk}
\usepackage{abstract}
\usepackage{xcolor}
\usepackage{graphicx}
\usepackage{booktabs}
\usepackage{array}
\usepackage{tabularx}
\usepackage{longtable}
\usepackage{enumitem}
\usepackage{listings}
\usepackage{hyperref}
\usepackage{url}
\usepackage{tikz}
\usetikzlibrary{arrows.meta, positioning, fit, backgrounds, shapes.geometric, calc}

% ---- arXiv-ish styling --------------------------------------------------------
\definecolor{accent}{HTML}{2E5BFF}
\definecolor{mute}{HTML}{606E7B}
\definecolor{soft}{HTML}{F5F6F8}
\definecolor{code}{HTML}{1F2933}

\hypersetup{
  colorlinks=true,
  linkcolor=accent,
  citecolor=accent,
  urlcolor=accent,
  pdftitle={Code Reviewer: an agentic LLM for pull-request review},
  pdfauthor={Team Messi}
}

\setlength{\droptitle}{-3em}
\titleformat*{\section}{\large\bfseries\sffamily}
\titleformat*{\subsection}{\normalsize\bfseries\sffamily}
\titleformat*{\subsubsection}{\normalsize\itshape\sffamily}

\renewcommand{\abstractnamefont}{\normalfont\bfseries\sffamily}
\renewcommand{\abstracttextfont}{\normalfont\small}

\lstdefinestyle{plain}{
  basicstyle=\ttfamily\footnotesize\color{code},
  backgroundcolor=\color{soft},
  frame=single,
  framerule=0pt,
  xleftmargin=8pt,
  xrightmargin=8pt,
  breaklines=true,
  showstringspaces=false,
  columns=fullflexible,
}
\lstset{style=plain}

% Allow flexible spacing so long unbreakable inline-code identifiers do not
% spill into the margin.
\sloppy
\emergencystretch=3em
\hyphenpenalty=200
\tolerance=2000

% Let \url break at underscores/dots/slashes/hyphens.
\makeatletter
\g@addto@macro\UrlBreaks{\do\_\do\.\do\/\do\-}
\makeatother

\newcommand{\code}[1]{\texttt{\small #1}}
\renewcommand{\path}[1]{\texttt{\small\color{mute} #1}}

% =============================================================================
\title{\textbf{\sffamily Code Reviewer}\\[2pt]
  \large\sffamily An agentic LLM that reviews pull requests on every push}

\author[1]{Cretu Cristian}
\author[1]{Cretu Luca}
\author[1]{Greholea Denis}
\author[1]{Gosa Bogdan}
\author[1]{Hiticas Paul}
\affil[1]{Team Messi}

\date{\today}

\begin{document}
\maketitle

\begin{abstract}
\noindent
We describe \emph{Code Reviewer}, an open-source pull-request review system that
runs as a single GitHub Action. The system combines a fine-tuned 9B-parameter
Qwen model, a per-repository retrieval-augmented generation (RAG) index, and a
\code{smolagents}-based tool loop into one pipeline that, on every PR open or
push, fetches the diff, indexes the head commit into a vector database,
investigates the change with a fixed set of tools, and submits one batched
review with inline comments and an explicit verdict
(\textsc{request\_changes}, \textsc{approve}, or \textsc{comment}). The model
is trained in two stages on top of \code{Qwen/Qwen3.5-9B}: a QLoRA supervised
finetune on the Microsoft CodeReviewer Comment Generation dataset, followed by
a GRPO~\cite{shao2024deepseekmath} stage with a Claude Haiku LLM-as-judge
reward and an explicit length penalty. A separate quantization pipeline
produces a 5.2\,GB \code{Q4\_K\_M} GGUF for local serving on Apple Silicon. We
document the full architecture, training recipes, retrieval design, agent
prompt, on-demand framework playbook system, and an end-to-end evaluation on
the CodeReviewer test split.
\end{abstract}

% ============================================================================
\section{Introduction}

Manual pull-request review is the de-facto quality gate of modern software
development, but it scales poorly: reviewers context-switch, comments are
inconsistent across people and repositories, and the cycle time blocks
merges. Generic LLM reviewers help in part but suffer from two recurring
failures: (1) they only see the diff and hallucinate APIs that do not exist in
the surrounding code, and (2) they cannot \emph{act} -- they cannot read
related files, run linters, or query git history.

This report documents an end-to-end alternative we built and shipped as a
GitHub Action. The system is structured as three loosely coupled components
(Figure~\ref{fig:overview}):
\begin{enumerate}[leftmargin=*]
  \item A \textbf{RAG sidecar} (\path{rag/}) that ingests the consumer
        repository on every run into a per-repository ChromaDB collection
        backed by \code{nomic-ai/CodeRankEmbed} embeddings.
  \item An \textbf{agent} (\path{agentic/}) built on
        \code{smolagents.CodeAgent} with a fixed toolbox for retrieval,
        execution, history, framework playbooks, and review actions.
  \item A two-stage \textbf{training pipeline}
        (\path{sft/}, \path{rlhf/}) that produces the policy model from
        \code{Qwen/Qwen3.5-9B}: a QLoRA SFT stage on the Microsoft
        CodeReviewer dataset followed by a GRPO~\cite{shao2024deepseekmath}
        stage with an LLM-as-judge reward.
\end{enumerate}

\begin{figure}[h]
\centering
\begin{tikzpicture}[
  node distance=8mm and 14mm,
  every node/.style={font=\footnotesize\sffamily},
  box/.style={draw=mute, fill=soft, rounded corners=2pt, minimum width=32mm,
              minimum height=8mm, align=center},
  acc/.style={draw=accent, fill=white, rounded corners=2pt, minimum width=32mm,
              minimum height=8mm, align=center, text=accent},
  arr/.style={-{Latex[length=2mm]}, draw=mute}]
  \node[box] (gh)  {GitHub PR event};
  \node[box, right=of gh] (ckt) {checkout HEAD};
  \node[box, right=of ckt] (rag) {RAG sidecar\\\scriptsize ingest};
  \node[acc, below=of rag] (agent) {agent loop\\\scriptsize tools $\le$ 20 calls};
  \node[box, left=of agent] (model) {fine-tuned\\Qwen3.5-9B};
  \node[box, left=of model] (out)   {one batched\\PR review};
  \draw[arr] (gh) -- (ckt);
  \draw[arr] (ckt) -- (rag);
  \draw[arr] (rag) -- (agent);
  \draw[arr] (agent) -- (model);
  \draw[arr] (model) -- (out);
\end{tikzpicture}
\caption{High-level flow. The Action checks out the PR head, starts the RAG
service, ingests the head into a Chroma collection, then runs the agent
loop. The agent buffers comments and emits one final review.}
\label{fig:overview}
\end{figure}

The remainder of this report walks through each component in the order the
code is organized in the repository.

% ============================================================================
\section{Repository Layout}

\begin{tabular}{@{}p{0.27\textwidth}p{0.69\textwidth}@{}}
\toprule
\textbf{Path} & \textbf{Purpose} \\
\midrule
\path{action.yml}           & Composite GitHub Action entry point. \\
\path{agentic/}             & Agent orchestration, tools, skills, review state. \\
\path{rag/}                 & FastAPI service for ingestion and semantic retrieval. \\
\path{sft/training/}        & QLoRA SFT trainer for \code{Qwen3.5-9B}. \\
\path{sft/data/}            & Dataset download, filtering, preprocessing, split. \\
\path{sft/eval/}            & Automated metrics, LLM-as-judge, defect metrics. \\
\path{rlhf/training/}       & GRPO trainer with a Claude Haiku judge reward. \\
\path{quantize/}            & llama.cpp-based pipeline producing a GGUF artifact. \\
\path{config/}              & Dynaconf configuration (dev / prod profiles). \\
\path{docs/}                & Operator runbooks (quantize, host on Mac). \\
\path{examples/review.yml}  & Consumer-repo workflow template. \\
\bottomrule
\end{tabular}

% ============================================================================
\section{The GitHub Action}
\label{sec:action}

The Action is a \emph{composite} action defined in \path{action.yml}. It
exposes seven inputs (\code{model-api-base}, \code{model-id},
\code{model-api-key}, \code{github-token}, \code{max-agent-steps} default
\code{20}, \code{comment-budget} default \code{10}, \code{python-version}
default \code{3.11}) and is intended to be invoked from a consumer workflow
such as \path{examples/review.yml}.

\paragraph{Steps executed inside the consumer runner.}
\begin{enumerate}[leftmargin=*,itemsep=2pt]
  \item \textbf{Validate event.} Refuses to run unless the trigger is
        \code{pull\_request} or \code{pull\_request\_target} and a PR number
        is present in the payload.
  \item \textbf{Install agent dependencies.} A pinned set is installed with
        CPU-only \code{torch} first to avoid pulling $\sim$2\,GB of CUDA
        wheels on the CPU GitHub runner; \code{transformers} is pinned to
        4.x because \code{nomic-ai/CodeRankEmbed}'s remote-code
        \code{NomicBertModel} rejects a keyword added in 5.x.
  \item \textbf{Start RAG sidecar.} Backgrounds \code{python -m rag.main}
        under \code{DYNACONF\_APP\_PROFILE=prod} and polls
        \code{http://localhost:8000/health} for up to two minutes.
  \item \textbf{Ingest target repository.}
        \code{python -m rag.scripts.ingest\_local} walks
        \code{\$GITHUB\_WORKSPACE} and posts files into the Chroma
        collection keyed by \code{\$REPO\_ID} (Section~\ref{sec:rag}).
  \item \textbf{Run review agent.} \code{python -m agentic.entrypoint}
        executes with the env vars wired from the Action inputs
        (Section~\ref{sec:agent}).
  \item \textbf{Upload artifact.} \code{review\_output.json} and
        \code{rag.log} are uploaded under the artifact name
        \code{ai-code-review} regardless of whether the agent succeeded.
\end{enumerate}

The Action is intentionally \emph{self-contained}: no Docker, no external
service besides whatever the consumer points \code{MODEL\_API\_BASE} at.

% ============================================================================
\section{Retrieval-Augmented Generation (RAG)}
\label{sec:rag}

The RAG service lives under \path{rag/} and exposes a small REST surface
backed by ChromaDB.

\subsection{Service shape}

\code{rag/main.py} constructs a single \code{FastAPI} app
(\code{code\_reviewer\_rag\_api}) and mounts the router from
\code{rag/routes.py}. Three endpoints are exposed:

\begin{itemize}[leftmargin=*]
  \item \code{GET /health} -- liveness probe used by the Action's polling
        loop.
  \item \code{POST /api/v1/ingest/\{repository\_id:path\}} -- accepts an
        \code{IngestionRequest} (a list of \code{SourceCodeSnippet} objects)
        and returns an \code{IngestionAcknowledgement}.
  \item \code{POST /api/v1/retrieve/\{repository\_id:path\}} -- accepts a
        \code{SearchQuery} (query text, \code{max\_results} clamped to
        $[1,20]$, optional metadata filter) and returns a
        \code{RetrievalResults} object containing matches with cosine
        similarity scores.
\end{itemize}

A \code{DELETE /api/v1/repo/\{id\}/file?path=...} endpoint removes all
snippets for a given file from the collection but is not used by the
default agent flow.

Schemas are defined as Pydantic models in \code{rag/schemas.py}. The
relevance score returned to clients is
$\mathrm{score} = 1 - \mathrm{distance}$, where \code{distance} is
ChromaDB's cosine distance.

\subsection{Per-repository collections}

\code{rag/database.py} wraps \code{chromadb.PersistentClient} with a
\code{RepositoryVectorDatabase} that creates one collection per repository.
Names are produced by \code{sanitize\_repository\_identifier}: non-alphanumeric
characters are replaced with hyphens, length is clamped to $[3, 63]$, and
leading/trailing characters are forced to alphanumeric -- matching ChromaDB's
collection-name constraint. Each collection is created with
\code{metadata=\{"hnsw:space": "cosine"\}} so similarity is cosine.

\subsection{Embeddings}

\code{rag/embeddings.py} implements
\code{TrustedSentenceTransformerEmbedding}, a Chroma \code{EmbeddingFunction}
that wraps \code{SentenceTransformer(model\_name, trust\_remote\_code=True)}.
The \code{trust\_remote\_code=True} flag is required because
\code{nomic-ai/CodeRankEmbed} ships a custom \code{NomicBertModel}
architecture. The same \code{embedding\_service} module-level instance is
used both at ingestion time and at query time, ensuring symmetric encoding.
The model identifier is read from \code{settings.EMBEDDINGS.MODEL\_NAME},
which is set to \code{nomic-ai/CodeRankEmbed} in both \code{config.dev.yaml}
and \code{config.prod.yaml}.

\subsection{Ingestion pipeline}

\code{rag/service.py} implements \code{ingest\_repository}. For each
snippet:

\begin{enumerate}[leftmargin=*]
  \item Base metadata is augmented with \code{file\_path},
        \code{repository\_id}, and \code{chunk\_index}; any
        \code{metadata\_overrides} from the request are merged on top.
  \item A deterministic ID is computed by
        \code{generate\_deterministic\_id} in \code{service.py} as the SHA-256
        of the string
        \code{\{repo\_id\}:\{file\_path\}:\{chunk\_index\}:sha256(content)}.
        This is what makes ingestion idempotent: re-running on the same
        content updates the existing entry instead of duplicating it.
  \item All snippets are sent through \code{collection.upsert(ids, documents,
        metadatas)} in a single call, so embedding happens in one batch.
\end{enumerate}

The CI ingester (\code{rag/scripts/ingest\_local.py}) walks the local
filesystem rather than fetching files from GitHub. It applies three filters
before sending snippets in batches of 25:

\begin{itemize}[leftmargin=*]
  \item Directory denylist:
        \code{\{.git, node\_modules, .venv, venv, \_\_pycache\_\_, .chroma\_data,
        dist, build, .idea, .github, .pytest\_cache, .mypy\_cache,
        .ruff\_cache\}}.
  \item Extension allowlist read from
        \code{settings.INGESTION.SUPPORTED\_EXTENSIONS}
        (\code{.py, .md, .ts, .js, .go, .java, .rs}).
  \item Per-file size cap \code{MAX\_FILE\_BYTES = 200{,}000}.
\end{itemize}

Each file becomes a single snippet with \code{chunk\_index=0} and language
metadata derived from its suffix. The script is the only one of the two
ingest scripts wired into the Action;
\code{rag/scripts/ingest\_repository.py} is an older variant that fetches
files via the GitHub raw API and is retained for ad-hoc indexing of remote
repositories.

\subsection{Query path}

\code{retrieve\_semantic\_context} delegates the actual search to
\code{collection.query(query\_texts=[q], n\_results=k, where=...,
include=["documents","metadatas","distances"])}. The list comprehension
that walks
\code{query\_results["documents"][0]} constructs \code{RetrievalMatch}
objects with
\code{relevance\_score = 1.0 - distance}. The retrieval HTTP path is hit
by exactly one consumer in production: the agent's
\code{SemanticSearchTool}.

% ============================================================================
\section{Configuration}

\code{config/config\_manager.py} reads \code{DYNACONF\_APP\_PROFILE} (must
be set; refusing to start otherwise) and constructs a
\code{Dynaconf} instance that layers:
\code{config/config.yaml}, \code{config/config.\{profile\}.yaml}, and
\code{config/.secrets.\{profile\}.yaml}. A small \code{SettingsWrapper}
raises if a queried attribute is missing rather than silently returning
\code{None}, which catches typos in the YAML keys at the point of use.

In the shipped Action the profile is \code{prod}; for local dry-runs the
README documents \code{DYNACONF\_APP\_PROFILE=dev}. The two YAML files are
identical at present; they exist as the seam for future
profile-specific overrides (e.g. a different embedding model in dev).

% ============================================================================
\section{The Agent}
\label{sec:agent}

\code{agentic/} contains the policy-side of the system: the entry point,
the agent constructor, the tool definitions, the review-state machine, and
the skills loader.

\subsection{Construction}

\code{agentic/agent.py} is twelve lines:

\begin{lstlisting}
from smolagents import CodeAgent, LiteLLMModel
from agentic.config import API_BASE, MODEL_ID
from agentic.tools import TOOLS

def build_model():
    return LiteLLMModel(model_id=MODEL_ID, api_base=API_BASE, max_tokens=4096)

def build_agent():
    return CodeAgent(tools=TOOLS, model=build_model())
\end{lstlisting}

\code{LiteLLMModel} routes through the LiteLLM SDK, which means any
OpenAI-compatible endpoint (vLLM, TGI, HF Inference Endpoints, mlx-lm,
\code{llama-server}, Together, Fireworks, etc.) can be substituted with no
code changes. \code{agentic/config.py} loads a local \path{agentic/.env}
file if one exists, defaults \code{API\_BASE} to \code{http://localhost:1234/v1}
(LM Studio), and \emph{prefixes the model id with} \code{openai/} so LiteLLM
treats the endpoint as OpenAI-compatible. It also seeds \code{OPENAI\_API\_KEY}
with \code{dummy} if the env var is missing, which is what allows
unauthenticated local endpoints to work.

\subsection{Entry point and prompt}
\label{sec:prompt}

\code{agentic/entrypoint.py} is the binary that the Action invokes. It reads
five required env vars (\code{GITHUB\_TOKEN}, \code{REPO\_ID}/%
\code{GITHUB\_REPOSITORY}, \code{PR\_NUMBER}, \code{HEAD\_SHA}, plus
\code{API\_BASE}), then:

\begin{enumerate}[leftmargin=*]
  \item Calls \code{GitHubClient.get\_pr}, \code{.get\_pr\_diff} (which uses
        the \code{application/vnd.github.v3.diff} Accept header), and
        \code{.get\_existing\_review\_comments}.
  \item Configures the singleton \code{REVIEW\_STATE}
        (Section~\ref{sec:review-state}) with the GitHub client, repo, PR
        number, head SHA, and comment budget.
  \item Truncates the diff to 60{,}000 characters via \code{\_truncate\_diff},
        splitting the budget in half between the head and tail with an
        explicit truncation marker.
  \item Runs \code{detect\_skills(diff=truncated\_diff)}
        (Section~\ref{sec:skills}) to produce a catalog of framework
        playbooks the agent \emph{can} load on demand.
  \item Assembles the task prompt by concatenating, in order: the
        \code{SYSTEM\_INSTRUCTIONS} string; the optional skills catalog;
        repo, PR, author, base$\leftarrow$head, file/line counts; the PR body;
        the truncated diff in a \code{diff}-fenced code block; and an
        \code{ALREADY FLAGGED} block listing the first 30 existing inline
        comments verbatim so the model does not re-emit them.
  \item Invokes \code{agent.run(task, max\_steps=MAX\_AGENT\_STEPS)} inside
        a \code{try/except} that, on any uncaught exception, sets the
        verdict to \textsc{comment} so the review still posts.
  \item Calls \code{REVIEW\_STATE.submit()} unconditionally.
\end{enumerate}

The \code{SYSTEM\_INSTRUCTIONS} string is the behavioural contract. It
encodes a three-step protocol: \emph{investigate} (require at least one
investigation tool call before any \code{post\_comment}), \emph{post
findings}, and \emph{verdict} (exactly one of the three verdict tools,
followed by \code{final\_answer("done")}). It also lists nine failure-mode
categories the reviewer should scan for: error-path state restoration, PII
in logs, hardcoded secrets/fallbacks, input validation, error UX,
auth/AuthZ gaps, crypto mistakes, injection, and framework gotchas.

\subsection{Tools}
\label{sec:tools}

All tools subclass \code{smolagents.Tool} and live under
\path{agentic/tools/}. Table~\ref{tab:tools} summarizes them; each is
described in detail below.

\begin{table}[h]
\centering
\small
\begin{tabular}{@{}p{0.25\textwidth}p{0.16\textwidth}p{0.50\textwidth}@{}}
\toprule
\textbf{Tool} & \textbf{Module} & \textbf{Purpose} \\
\midrule
\code{semantic\_search}       & retrieval & RAG over the ingested codebase. \\
\code{search\_keyword}        & retrieval & \code{grep -rn} across the checkout. \\
\code{search\_symbol}         & retrieval & AST walk for Python definitions. \\
\code{get\_file}              & retrieval & Read a file or a \code{start-end} line range. \\
\code{run\_tests}             & execution & \code{pytest} on the working tree. \\
\code{run\_linter}            & execution & \code{ruff check} on a path or repo. \\
\code{run\_typecheck}         & execution & \code{mypy} on a path or repo. \\
\code{run\_snippet}           & execution & Sandboxed Python exec with a 5\,s timeout. \\
\code{check\_history}         & history   & \code{git blame} on a path or line range. \\
\code{get\_pr\_metadata}      & history   & GitHub REST \code{GET /pulls/\{n\}}; falls back to \code{git log -1}. \\
\code{get\_team\_conventions} & history   & Read \code{.cursorrules}, \code{CONTRIBUTING.md}, \code{CLAUDE.md}. \\
\code{load\_skill}            & skills    & Pull a framework playbook into context. \\
\code{post\_comment}          & action    & Buffer one inline comment on \code{file:line}. \\
\code{propose\_patch}         & action    & Buffer a GitHub-style \code{suggestion} block. \\
\code{request\_changes}       & action    & Verdict: \textsc{request\_changes}. \\
\code{approve}                & action    & Verdict: \textsc{approve}. \\
\code{comment\_only}          & action    & Verdict: \textsc{comment}. \\
\bottomrule
\end{tabular}
\caption{The agent's tool surface. All tools are passed to \code{CodeAgent}
in the order listed in \path{agentic/tools/\_\_init\_\_.py}.}
\label{tab:tools}
\end{table}

\subsubsection{Retrieval tools}

\code{SemanticSearchTool} POSTs to
\code{\{RAG\_URL\}/api/v1/retrieve/\{REPO\_ID\}} with the agent's query and
clamps \code{k} to $[1,20]$. \code{SearchKeywordTool} shells out to
\code{grep -rn} from the repository root, with an optional
\code{--include path\_glob}. \code{SearchSymbolTool} performs a recursive
\code{os.walk} over the workspace, parses every \code{.py} file with the
standard \code{ast} module, and walks the tree for
\code{FunctionDef}, \code{AsyncFunctionDef}, or \code{ClassDef} nodes
matching \code{name}. \code{GetFileTool} reads a path with optional
\code{"start-end"} line slicing.

\subsubsection{Execution tools}

All four execution tools shell out to subprocesses with
\code{capture\_output=True}. \code{RunSnippetTool} writes the snippet to
\code{/tmp/snippet.py} and runs \code{python3} with a five-second timeout;
the temp file is cleaned up in a \code{finally} block. Only
\code{language in \{"python","python3"\}} is accepted.

\subsubsection{History and metadata}

\code{CheckHistoryTool} invokes \code{git blame [-L line\_range] path}.
\code{GetPrMetadataTool} prefers the GitHub API path when
\code{GITHUB\_TOKEN} is set and falls back to \code{git log -1} otherwise,
so the same tool works locally and in CI. \code{GetTeamConventionsTool}
reads up to four files at known paths (\code{.cursorrules},
\code{.github/CONTRIBUTING.md}, \code{CONTRIBUTING.md}, \code{CLAUDE.md})
and returns their concatenation; this is the agent's hook for
respecting team-specific style.

\subsubsection{Action tools and review state}
\label{sec:review-state}

The action tools are the only ones that mutate the review. They all
delegate to the module-level singleton \code{REVIEW\_STATE} defined in
\code{agentic/review\_state.py}. The state machine enforces four
invariants:

\begin{enumerate}[leftmargin=*]
  \item \textbf{Dedupe by \code{(file, line)}.} \code{add\_comment} returns
        \code{"duplicate"} if a comment already exists on that pair; the
        model has a tendency to re-emit findings across steps.
  \item \textbf{Hard comment budget.} Returns \code{"budget\_exhausted"}
        once \code{len(comments) >= comment\_budget}.
  \item \textbf{Comments frozen after verdict.} Once
        \code{set\_verdict} succeeds, \code{add\_comment} returns
        \code{"finalized"}.
  \item \textbf{First-verdict-wins.} \code{set\_verdict} returns
        \code{False} on the second call.
\end{enumerate}

\code{submit()} POSTs a single
\code{POST /repos/\{repo\}/pulls/\{pr\}/reviews} payload with the buffered
\code{comments}, the chosen \code{event}, and a short body summarizing the
number of findings. The submission has \emph{two} fallback paths:

\begin{itemize}[leftmargin=*]
  \item If GitHub returns \code{422} on an \textsc{approve} verdict (the
        default \code{GITHUB\_TOKEN} cannot self-approve), the code
        downgrades the event to \textsc{comment} and retries while keeping
        all inline comments.
  \item If either submission still fails (e.g. an inline comment refers to
        a line that does not exist on the diff), the code folds the
        per-comment bodies into the review body as a Markdown list and
        retries with an empty \code{comments} array; failing that, it
        dumps the payload to \code{review\_output.json} and to stdout.
\end{itemize}

When the agent runs without a configured GitHub client (local dry-run),
\code{submit()} skips the network call entirely and only writes
\code{review\_output.json}.

\subsection{Skills: framework playbooks on demand}
\label{sec:skills}

\code{agentic/skills/} contains five markdown playbooks
(\code{async-js.md}, \code{nextjs.md}, \code{react.md}, \code{supabase.md},
\code{vite.md}), each with YAML-style frontmatter declaring its name,
description, and \code{triggers}. Trigger fields are:
\code{package\_json\_dep}, \code{package\_json\_dep\_prefix},
\code{pyproject\_dep}, \code{cargo\_dep}, \code{go\_mod\_module},
\code{files}, and \code{diff\_extensions}. A skill matches if \emph{any}
trigger matches (OR semantics).

The two-stage design in \code{agentic/skills/loader.py} keeps the cost on
the prompt small:

\begin{enumerate}[leftmargin=*]
  \item \textbf{Catalog} (\code{detect\_skills}). At task setup, the loader
        parses the consumer repository's \code{package.json},
        \code{pyproject.toml}, \code{Cargo.toml}, and \code{go.mod}, plus
        the set of file extensions present in the diff, and selects all
        skills whose triggers match. Each entry costs roughly one line on
        the prompt (\code{name} + description). The catalog is injected
        under the heading \code{AVAILABLE PLAYBOOKS}.
  \item \textbf{Body} (\code{load\_skill\_body}). The
        \code{LoadSkillTool} fetches the full body of a named skill on
        demand and caches the \emph{name} (not the body) on the instance
        so the agent receives a polite no-op on the second call instead
        of duplicating the body in the conversation.
\end{enumerate}

The frontmatter parser is intentionally hand-rolled rather than pulling in
PyYAML for this one use case; it handles only top-level scalar keys and
the \code{triggers:} block (lists or scalar values). The package.json
parser handles \code{dependencies}, \code{devDependencies}, and
\code{peerDependencies}. The \code{pyproject.toml} parser uses two regexes
to catch both Poetry-style (\code{name = "1.0"}) and PEP-621
(\code{"name>=1.0"}) dependency declarations without importing a TOML
parser.

% ============================================================================
\section{Model Training}

The policy model is \code{Qwen/Qwen3.5-9B}, a 9B-parameter decoder-only
causal LM from Alibaba's Qwen team. We train in two stages.

\subsection{Stage 1: Supervised fine-tuning (QLoRA)}

\paragraph{Data pipeline.} The pipeline is split into four numbered
modules under \path{sft/data/}:

\begin{enumerate}[leftmargin=*]
  \item \code{download.py} fetches the Microsoft \emph{CodeReviewer
        Comment Generation} dataset (846\,MB) from Zenodo record 6900648
        and extracts it under \path{data/raw/}.
  \item \code{filter.py} removes obviously non-actionable comments before
        training (training split only; valid/test pass through). The
        filters are: length below 10 characters, exact-match against a
        38-entry noisy lexicon (\code{lgtm}, \code{done}, \code{ok},
        \code{+1}, \code{wip}, \code{thanks}, $\ldots$), and seven regex
        patterns covering bare URLs, ``see \#N'', ``duplicate of \#N'',
        and punctuation-only comments.
  \item \code{preprocess.py} converts each surviving \code{(patch, msg)}
        pair into a three-message chat conversation: a fixed
        \code{SYSTEM\_PROMPT} (the SFT system prompt; reproduced at the
        end of this section), a user message wrapping the normalized diff
        in a fenced \code{diff} block, and the assistant message with the
        ground-truth comment. The patch normalizer rewrites
        CodeReviewer's \code{<add>}/\code{<del>} tags to the more
        common \code{+}/\code{-} prefixes.
  \item \code{split.py} shuffles the preprocessed train split with seed
        42 and emits an 80/10/10 train/val/test partition into
        \path{data/processed/\{train,val,test\}.jsonl}.
\end{enumerate}

\paragraph{Training loop.} \code{sft/training/sft.py} imports
\code{unsloth} \emph{before} \code{trl}/\code{transformers} so the patch
that powers Unsloth's accelerated kernels takes effect. The script reads
\path{sft/training/config.yaml} and loads the model via
\code{FastLanguageModel.from\_pretrained} with \code{max\_seq\_length=2048},
\code{load\_in\_4bit=True}, and the auto dtype (\code{bf16} on A100). LoRA
adapters are then attached with:

\begin{tabular}{@{}ll@{}}
rank $r$              & 16 \\
$\alpha$              & 16 \\
dropout               & 0  \\
target modules        & \code{q\_proj, k\_proj, v\_proj, o\_proj, gate\_proj, up\_proj, down\_proj} \\
\end{tabular}

\noindent
Training uses TRL's \code{SFTTrainer} with
\code{per\_device\_train\_batch\_size=32},
\code{gradient\_accumulation\_steps=1}, \code{learning\_rate=2e-4},
\code{lr\_scheduler="cosine"}, \code{weight\_decay=0.01},
\code{max\_grad\_norm=1.0}, \code{bf16=true}, \code{optim="adamw\_8bit"}, and
\code{num\_train\_epochs=1}. The chat-format conversion is done at
\code{Dataset.map} time via \code{tokenizer.apply\_chat\_template(...,
add\_generation\_prompt=False)}, so the model is trained on the full
assistant turn including the response. The script supports two debug knobs:
\code{--small-run} (1K/100 train/val examples) and \code{--max-examples N}.

The resulting adapter is published as
\href{https://huggingface.co/cristicretu/code-reviewer-lora}{\code{cristicretu/code-reviewer-lora}}.

\subsection{Stage 2: RLHF with GRPO and a Claude judge}

\code{rlhf/training/grpo.py} starts from the SFT adapter and continues with
Group Relative Policy Optimization~\cite{shao2024deepseekmath} via TRL's
\code{GRPOTrainer}. For each prompt the trainer samples
\code{num\_generations=4} completions at \code{temperature=0.8} with
\code{max\_prompt\_length=512} and \code{max\_completion\_length=512}; the
group-relative advantage replaces the value head used in vanilla PPO.

\paragraph{Reward.} The reward is an LLM-as-judge call to
\code{claude-haiku-4-5-20251001} through the Anthropic SDK. The judge
prompt (\code{JUDGE\_PROMPT} at the top of the file) asks for a single
integer 1--5 with the rubric:

\begin{itemize}[leftmargin=*,nosep]
  \item \textbf{5} -- real issue, exact location, explanation, concrete fix.
  \item \textbf{4} -- real issue but vague about location, cause, or fix.
  \item \textbf{3} -- partially correct.
  \item \textbf{2} -- generic observation applicable to any diff.
  \item \textbf{1} -- factually wrong, refers to absent code, or toxic.
\end{itemize}

After parsing the integer, the reward function subtracts a length penalty
of $\min(0.5, \max(0, (w-100)/100) \cdot 0.5)$ where $w$ is the completion
word count. This caps verbosity at no extra cost for $\le 100$ words and
fades the score down by up to 0.5 for completions over 200 words.

Robustness is layered in three places:

\begin{itemize}[leftmargin=*]
  \item \emph{API errors.} \code{anthropic.RateLimitError} retries up to
        four times with $2^a \cdot 5$ second backoff;
        \code{APITimeoutError}, \code{APIConnectionError}, and
        \code{InternalServerError} retry with $2^a \cdot 2$ second backoff;
        repeated failure returns the floor score \code{1.0}.
  \item \emph{Concurrency.} A module-level
        \code{ThreadPoolExecutor(max\_workers=8)} fans out judge calls
        across the group of 4 completions $\times$ batch.
  \item \emph{Parse errors.} A non-integer judge response (e.g. the model
        emitting a \code{ToolUse} block instead of \code{TextBlock}) also
        floors to 1.0.
\end{itemize}

\paragraph{Optimizer.} Training runs one epoch on the first 250 prompts of
the preprocessed train split, at \code{learning\_rate=2e-6},
\code{per\_device\_train\_batch\_size=1},
\code{gradient\_accumulation\_steps=4}, \code{bf16=true}. The model is
loaded from the SFT adapter with \code{load\_in\_4bit=True}. After
training, \code{model.save\_pretrained(f"\{output\_dir\}/final")} writes a
PEFT-format checkpoint and, if \code{GRPO\_HUB\_MODEL\_ID} (or the YAML
\code{hub\_model\_id}) is set, pushes to the Hub. The published checkpoint
is \href{https://huggingface.co/cretu-luca/code-reviewer-grpo}%
{\code{cretu-luca/code-reviewer-grpo}}.

% ============================================================================
\section{Quantization for Apple Silicon}

\code{quantize/quantize.py} produces a 4-bit GGUF artifact suitable for
\code{llama-server} on a 24\,GB M4 Pro. The pipeline composes four
upstream tools from \code{ggml-org/llama.cpp}:

\begin{enumerate}[leftmargin=*]
  \item \code{convert\_hf\_to\_gguf.py base\_dir \texttt{-{}-}outtype f16
        \texttt{-{}-}outfile base-f16.gguf}.
  \item \code{convert\_lora\_to\_gguf.py adapter\_dir \texttt{-{}-}base
        base\_dir \texttt{-{}-}outfile grpo-lora.gguf}.
  \item \code{llama-export-lora -m base-f16.gguf \texttt{-{}-}lora
        grpo-lora.gguf -o fused-f16.gguf}.
  \item \code{llama-quantize fused-f16.gguf out.gguf Q4\_K\_M}.
\end{enumerate}

The non-trivial step the script handles itself is \emph{staging}: it
materializes a directory whose weights/\code{config.json} come from the
base model snapshot but whose tokenizer files
(\code{tokenizer.json}, \code{tokenizer\_config.json},
\code{chat\_template.jinja}, \code{special\_tokens\_map.json},
\code{vocab.json}, \code{merges.txt}) come from the GRPO adapter
repository. This ensures the GGUF metadata embeds the chat template the
model was \emph{trained} with, not the base's default. Without this, the
served model drifts off-distribution and can stall inside the thinking
block.

\code{which\_or\_die} resolves the two binaries (\code{llama-quantize},
\code{llama-export-lora}); the latter is not in the Homebrew formula and
must be built from source, so the function also searches
\path{./.quantize-work/llama.cpp-src/build/bin}.

The default quantization \code{Q4\_K\_M} produces a $\sim$5.2\,GB GGUF
that resides in $\sim$6.5\,GB of RAM at inference, leaving comfortable
headroom on a 24\,GB Mac for KV cache and the rest of the system. The
choice is justified in \path{docs/quantize-on-mac.md} with a per-step
walltime/RAM/disk table on M4 Pro: $\sim$10 minutes wall clock total,
peaking at $\sim$40\,GB free disk during \code{llama-export-lora}. The
published GGUF lives at
\href{https://huggingface.co/cretu-luca/code-reviewer-grpo-GGUF}%
{\code{cretu-luca/code-reviewer-grpo-GGUF}}.

% ============================================================================
\section{Evaluation}

Evaluation is split between fast automated metrics
(\path{sft/eval/metrics.py}, \path{sft/eval/defect\_metrics.py}) and
LLM-as-judge runs (\path{sft/eval/judge.py},
\path{sft/eval/deepeval\_judge.py}).

\subsection{Automated metrics}

\code{compute\_metrics(preds, refs)} returns three reference-based scores:

\begin{itemize}[leftmargin=*,nosep]
  \item \textbf{ChrF} via \code{sacrebleu.corpus\_chrf}.
  \item \textbf{ROUGE-L} F1 via \code{rouge\_score.rouge\_scorer.RougeScorer(["rougeL"], use\_stemmer=True)}.
  \item \textbf{CodeBERTScore} F1 from the \code{code\_bert\_score} package,
        falling back to \code{None} if the optional dependency is absent.
\end{itemize}

BLEU is intentionally omitted (corpus BLEU on review comments scores
below 10 and is uninformative).

\code{defect\_metrics.py} adds four reference-free or
diff-relative metrics, all token/regex based and fast enough to run on
an M1 Pro in seconds:

\begin{itemize}[leftmargin=*,nosep]
  \item \textbf{Hallucination rate.} For each prediction, extracts
        identifiers via \code{[a-zA-Z\_][a-zA-Z0-9\_]\{2,\}}, removes a
        50-word English stoplist, and flags a prediction as hallucinated
        if more than 30\% of its identifiers do not appear in the diff.
  \item \textbf{Defect F1.} Approximates issue overlap by computing
        identifier-set precision/recall against the reference comment
        and reporting their harmonic mean.
  \item \textbf{Coherence.} A weighted score
        (\code{0.25 sent + 0.25 len + 0.20 code-ratio + 0.30 structure})
        rewarding 3--8 sentences, average sentence length near 15 words,
        $\sim$15\% code tokens, and the presence of explanatory phrases
        (\textit{the issue is}, \textit{this could cause},
        \textit{to fix}, \textit{because}).
  \item \textbf{Toxicity rate.} A boolean match against five regex
        patterns covering insults, sarcasm, and gendered pronouns in
        evaluative position.
\end{itemize}

\subsection{LLM-as-judge}

\code{sft/eval/judge.py} scores three dimensions (accuracy, helpfulness,
specificity) on a 1--5 scale and supports two backends:

\begin{itemize}[leftmargin=*,nosep]
  \item \code{--backend gemini} (\code{gemini-2.5-flash}) with explicit
        rate-limit handling: a 13-second sleep between requests to respect
        the free-tier 5\,RPM cap, plus exponential backoff (up to 5
        attempts) on 429/503.
  \item \code{--backend anthropic} (\code{claude-sonnet-4-20250514}).
\end{itemize}

A separate \code{deepeval\_judge.py} adds G-Eval rubrics for
\emph{correctness}, \emph{specificity}, and \emph{actionability}, plus
the built-in \code{AnswerRelevancyMetric}, scored by a local Ollama model
configured via \code{LOCAL\_MODEL\_*} env vars.

\subsection{Headline results}

We compare three checkpoints on the CodeReviewer test split: the
\code{qwen3:8b} base, the SFT adapter
(\code{cristicretu/code-reviewer-lora}), and the GRPO model
(\code{cretu-luca/code-reviewer-grpo}). Automated metrics are computed on
100 examples; LLM-as-judge on 10 examples per model with Gemini 2.5 Flash
as the judge.

\begin{table}[h]
\centering
\small
\begin{tabular}{@{}lccc@{}}
\toprule
\textbf{Metric}      & \textbf{Base (qwen3:8b)} & \textbf{SFT}    & \textbf{RLHF}    \\
\midrule
Defect F1            & 0.001                    & 0.070           & \textbf{0.072}  \\
CodeBERTScore        & 0.000                    & 0.668           & \textbf{0.669}  \\
ChrF                 & 0.29                     & 13.45           & \textbf{13.93}  \\
ROUGE-L              & 0.001                    & 0.049           & \textbf{0.050}  \\
Coherence            & 0.009                    & 0.491           & \textbf{0.507}  \\
Hallucination rate   & 100\%                    & 76\%            & 81\%            \\
Toxicity rate        & 0\%                      & 0\%             & 0\%             \\
\bottomrule
\end{tabular}
\caption{Automated metrics, $n=100$. Both fine-tuned variants dwarf the
base on every reference-based metric; the base produces essentially empty
reviews (mean response length $\approx 3$ characters).}
\label{tab:auto}
\end{table}

\begin{table}[h]
\centering
\small
\begin{tabular}{@{}lcc@{}}
\toprule
\textbf{Dimension}   & \textbf{SFT ($n=10$)} & \textbf{RLHF ($n=8$)} \\
\midrule
Accuracy             & \textbf{2.40}/5.0     & 1.50/5.0              \\
Helpfulness          & \textbf{1.90}/5.0     & 1.25/5.0              \\
Specificity          & \textbf{3.70}/5.0     & 2.62/5.0              \\
\bottomrule
\end{tabular}
\caption{LLM-as-judge (Gemini 2.5 Flash). The independent judge prefers the
SFT model on all three dimensions, suggesting the GRPO stage over-optimized
the Claude Haiku judge it was trained against.}
\label{tab:judge}
\end{table}

\paragraph{Discussion.} Three findings stand out (cf.\ \path{EVAL\_RESULTS.md}):
(1) both fine-tunes dramatically outperform the base on every metric, which
is unsurprising since the base produces near-empty reviews;
(2) automated metrics show a small consistent edge for RLHF over SFT
(Defect F1 0.072 vs.\ 0.070, ChrF 13.93 vs.\ 13.45); and
(3) the cross-judge LLM evaluation reverses that ordering. Together these
suggest the GRPO stage learned to maximize the Claude Haiku reward signal
without an unambiguous gain in subjective review quality. The hallucination
metric (76--81\%) reads alarming in isolation but is a heuristic: any new
identifier introduced by a suggested fix counts against the score, which is
exactly what good reviews do.

\subsection{End-to-end agent evaluation}

The README also reports an end-to-end evaluation of the
\code{code-reviewer-4-bit} GGUF inside the \code{smolagents.CodeAgent}
pipeline on 15 examples from the CodeReviewer test set, with the local
\code{gemma-4-e4b-it} as judge. The three metrics are:

\begin{itemize}[leftmargin=*,nosep]
  \item \textbf{Task Completion} (binary; was a valid verdict set before
        the step budget): 13/15 = 86.7\%, mean 0.85.
  \item \textbf{Task Correctness} (G-Eval over the GRPO reward rubric):
        13/15 pass at $\ge 0.50$ threshold, mean 0.85.
  \item \textbf{Step Efficiency} (linear decay over the 10-call cap,
        defined in the README): 12/15 = 80\%, mean 0.82.
\end{itemize}

The two budget-exhausted runs share a failure mode: the model emitted
prose review text rather than Python tool calls, leaving \code{smolagents}
with nothing executable.

% ============================================================================
\section{Operations}

\subsection{Consumer integration}

Three GitHub repository secrets must be set: \code{MODEL\_API\_BASE},
\code{MODEL\_ID}, \code{MODEL\_API\_KEY}. \path{examples/review.yml}
is the recommended workflow:

\begin{lstlisting}
name: AI Code Review
on: { pull_request: { types: [opened, synchronize, reopened] } }
permissions: { contents: read, pull-requests: write }
concurrency: { group: ai-review-${{ github.event.pull_request.number }},
               cancel-in-progress: true }
jobs:
  review:
    runs-on: ubuntu-latest
    timeout-minutes: 25
    steps:
      - uses: actions/checkout@v4
        with: { ref: ${{ github.event.pull_request.head.sha }}, fetch-depth: 0 }
      - uses: cristicretu/code-reviewer@v1
        with:
          model-api-base: ${{ secrets.MODEL_API_BASE }}
          model-id:       ${{ secrets.MODEL_ID }}
          model-api-key:  ${{ secrets.MODEL_API_KEY }}
\end{lstlisting}

The \code{concurrency} block cancels in-flight reviews when a new commit
lands, which avoids paying for a stale review that the next push would
invalidate anyway.

\subsection{Self-hosting on Apple Silicon}

Two paths are documented in \path{docs/}. \path{quantize-on-mac.md}
covers building the \code{Q4\_K\_M} GGUF and serving it via
\code{llama-server -m \ldots --port 8080 --jinja -c 4096 -ngl 999}, which
exposes an OpenAI-compatible \code{/v1/chat/completions} endpoint with
\code{reasoning\_content} support. \path{host-on-mac.md} covers the
zero-cost demo path: convert with \code{mlx-lm} 4-bit, serve on
\code{127.0.0.1:8001}, then expose it via
\code{cloudflared tunnel --url http://127.0.0.1:8001} and paste the
resulting \code{trycloudflare.com} URL into the consumer repo's
\code{MODEL\_API\_BASE} secret.

% ============================================================================
\section{Failure Modes and Mitigations}

The behavior of the system on adversarial inputs is shaped largely by the
review-state machine and the agent's prompt:

\begin{itemize}[leftmargin=*]
  \item \textbf{Spurious comments.} The system prompt mandates at least one
        investigation tool call before any \code{post\_comment}. A run that
        skips this is, by the prompt's own definition, a failed review.
  \item \textbf{Comment flood.} The hard \code{comment\_budget} (default
        10) is enforced server-side in \code{ReviewState.add\_comment};
        the prompt also asks the agent to ``pick the highest-impact issues
        first.''
  \item \textbf{Duplicate findings on consecutive steps.} Dedupe by
        \code{(file, line)} in \code{add\_comment}.
  \item \textbf{Indecision.} \code{set\_verdict} is first-call-wins; the
        per-tool docstring tells the agent that subsequent verdict calls
        are no-ops.
  \item \textbf{Crashes mid-run.} \code{entrypoint.main} catches any agent
        exception, sets the verdict to \code{COMMENT}, and still posts
        \code{REVIEW\_STATE.submit()}.
  \item \textbf{\textsc{approve} rejected (422).} \code{ReviewState.submit}
        retries the same payload as \textsc{comment}; if that fails too,
        comments are folded into the body and the inline payload is
        emptied.
  \item \textbf{Existing comments.} The prompt lists up to 30 previously
        flagged comments under \code{ALREADY FLAGGED} so the agent does
        not re-emit them on a subsequent push.
\end{itemize}

% ============================================================================
\section{Limitations and Future Work}

The system is small enough to enumerate its rough edges:

\begin{itemize}[leftmargin=*]
  \item \textbf{\code{search\_symbol} is Python-only.} It walks
        \code{.py} files with the standard \code{ast} module; TypeScript,
        Go, and Rust definitions are visible only through
        \code{search\_keyword}.
  \item \textbf{No execution sandbox.} \code{run\_tests},
        \code{run\_linter}, \code{run\_typecheck}, and \code{run\_snippet}
        all execute against the unsandboxed runner working directory. The
        snippet tool has only a five-second timeout. The Action should
        not be run on \code{pull\_request\_target} from forks without
        further isolation.
  \item \textbf{GRPO reward is single-judge.} The reward depends on one
        Claude Haiku call per completion. Table~\ref{tab:judge} shows the
        usual symptom: the cross-judge evaluation does not confirm the
        gain over SFT. A multi-judge or human-in-the-loop reward would
        likely help.
  \item \textbf{Skills catalog is short.} Five playbooks cover React,
        Next.js, Vite, Supabase, and async JS/TS. Python frameworks
        (Django, FastAPI), Go, Rust, and infra-as-code are not yet
        represented; the skill format is intentionally one file per
        framework so adding one is a contribution without code changes.
\end{itemize}

% ============================================================================
\section{Conclusion}

We have described a complete, deployable system for automated pull-request
review. The architecture is deliberately simple: a small RAG service, an
agent with seventeen tools (Table~\ref{tab:tools}), and a singleton state
machine that buffers and submits one batched review per run. Most of the
engineering is in the contract -- the agent's prompt, the tool surface, and
the review-state invariants -- rather than in the model. The trained
checkpoint is published as
\code{cretu-luca/code-reviewer-grpo} (and \code{...-GGUF} for Apple Silicon),
and the Action is consumable from any repository with three secrets and a
nine-line workflow.

% ============================================================================
\bibliographystyle{plain}
\begin{thebibliography}{9}

\bibitem{shao2024deepseekmath}
Shao, Z., Wang, P., Zhu, Q., Xu, R., Song, J., Bi, X., Zhang, H., Zhang, M.,
Li, Y., Wu, Y., Guo, D.
\newblock DeepSeekMath: Pushing the Limits of Mathematical Reasoning in Open
Language Models.
\newblock \emph{arXiv preprint arXiv:2402.03300}, 2024.
\newblock \url{https://arxiv.org/abs/2402.03300}.

\bibitem{li2022codereviewer}
Li, Z., Lu, S., Guo, D., Duan, N., Jannu, S., Jenks, G., Majumder, D.,
Green, J., Svyatkovskiy, A., Fu, S., Sundaresan, N.
\newblock Automating Code Review Activities by Large-Scale Pre-Training.
\newblock \emph{ESEC/FSE}, 2022.
\newblock Dataset: Zenodo record 6900648.

\bibitem{nomic2024coderankembed}
Nomic AI.
\newblock CodeRankEmbed: A high-quality code embedding model.
\newblock \url{https://huggingface.co/nomic-ai/CodeRankEmbed}.

\bibitem{qwen2024qwen35}
Qwen Team, Alibaba.
\newblock Qwen3.5 technical report.
\newblock \url{https://huggingface.co/Qwen/Qwen3.5-9B}.

\bibitem{unsloth}
Han, D., Han, M.
\newblock Unsloth: 2--5x faster LLM training with less memory.
\newblock \url{https://github.com/unslothai/unsloth}.

\bibitem{trl}
von Werra, L., et al.
\newblock TRL: Transformer Reinforcement Learning.
\newblock \url{https://github.com/huggingface/trl}.

\bibitem{smolagents}
Hugging Face.
\newblock \texttt{smolagents}: a barebones library for agents.
\newblock \url{https://github.com/huggingface/smolagents}.

\bibitem{chromadb}
Chroma.
\newblock The AI-native open-source embedding database.
\newblock \url{https://www.trychroma.com}.

\bibitem{llamacpp}
Gerganov, G., et al.
\newblock llama.cpp: LLM inference in C/C++.
\newblock \url{https://github.com/ggerganov/llama.cpp}.

\end{thebibliography}

\end{document}
